\def\tituloconcepto{La esfera de Bloch}
% Si el capitulo se compila solo (sin main.tex), esta cabecera abre el documento.
% Antes de % Si el capitulo se compila solo (sin main.tex), esta cabecera abre el documento.
% Antes de % Si el capitulo se compila solo (sin main.tex), esta cabecera abre el documento.
% Antes de \input{cabecera} cada capitulo define \tituloconcepto.
\ifdefined\enmain\else
  \documentclass[11pt,a4paper]{article}
  \input{preambulo}
  \begin{document}
  \title{\tituloconcepto}
  \author{\autor}
  \date{\today}
  \maketitle
\fi
 cada capitulo define \tituloconcepto.
\ifdefined\enmain\else
  \documentclass[11pt,a4paper]{article}
  % Preambulo compartido por todos los documentos. Compilar con XeLaTeX/LuaLaTeX
% (o Tectonic). Si se usa Overleaf, elegir el compilador XeLaTeX.
\usepackage[spanish,es-noshorthands]{babel}
\usepackage{amsmath,amssymb}
\usepackage{booktabs,array}
\usepackage[a4paper,margin=2.5cm]{geometry}
\usepackage{xcolor}
\usepackage{listings}
\usepackage[most]{tcolorbox}
\usepackage{quantikz}
\usepackage[hidelinks]{hyperref}

% ---- Autor (editar a gusto) -------------------------------------------------
\newcommand{\autor}{Genesis Silesky}

% ---- Notacion de Dirac y atajos ---------------------------------------------
% quantikz ya define \ket y \bra; DeclareRobustCommand los redefine sin error.
\DeclareRobustCommand{\ket}[1]{\left|#1\right\rangle}
\DeclareRobustCommand{\bra}[1]{\left\langle#1\right|}
\DeclareRobustCommand{\braket}[2]{\left\langle#1\middle|#2\right\rangle}
\newcommand{\C}{\mathbb{C}}
\newcommand{\R}{\mathbb{R}}
\DeclareMathOperator{\Tr}{Tr}

% ---- Cajas de apoyo --------------------------------------------------------
% "concepto": definiciones previas que hacen falta para entender el capitulo.
\newtcolorbox{concepto}[1]{breakable,colback=blue!4,colframe=blue!55!black,
  fonttitle=\bfseries,title={Concepto previo: #1}}
% "nota": aclaraciones y advertencias.
\newtcolorbox{nota}[1]{breakable,colback=orange!6,colframe=orange!70!black,
  fonttitle=\bfseries,title={#1}}

% ---- Estilos de listados ----------------------------------------------------
\lstdefinestyle{py}{language=Python,basicstyle=\ttfamily\footnotesize,
  breaklines=true,postbreak=\mbox{\textcolor{gray}{$\hookrightarrow$}\space},
  numbers=left,numberstyle=\tiny\color{gray},frame=single,
  commentstyle=\color{green!40!black}\itshape,
  keywordstyle=\color{blue!70!black}\bfseries,
  stringstyle=\color{red!60!black},showstringspaces=false,tabsize=4,
  columns=fullflexible,keepspaces=true}
\lstdefinestyle{salida}{basicstyle=\ttfamily\footnotesize,breaklines=true,
  frame=leftline,backgroundcolor=\color{gray!8},columns=fullflexible,
  keepspaces=true}

  \begin{document}
  \title{\tituloconcepto}
  \author{\autor}
  \date{\today}
  \maketitle
\fi
 cada capitulo define \tituloconcepto.
\ifdefined\enmain\else
  \documentclass[11pt,a4paper]{article}
  % Preambulo compartido por todos los documentos. Compilar con XeLaTeX/LuaLaTeX
% (o Tectonic). Si se usa Overleaf, elegir el compilador XeLaTeX.
\usepackage[spanish,es-noshorthands]{babel}
\usepackage{amsmath,amssymb}
\usepackage{booktabs,array}
\usepackage[a4paper,margin=2.5cm]{geometry}
\usepackage{xcolor}
\usepackage{listings}
\usepackage[most]{tcolorbox}
\usepackage{quantikz}
\usepackage[hidelinks]{hyperref}

% ---- Autor (editar a gusto) -------------------------------------------------
\newcommand{\autor}{Genesis Silesky}

% ---- Notacion de Dirac y atajos ---------------------------------------------
% quantikz ya define \ket y \bra; DeclareRobustCommand los redefine sin error.
\DeclareRobustCommand{\ket}[1]{\left|#1\right\rangle}
\DeclareRobustCommand{\bra}[1]{\left\langle#1\right|}
\DeclareRobustCommand{\braket}[2]{\left\langle#1\middle|#2\right\rangle}
\newcommand{\C}{\mathbb{C}}
\newcommand{\R}{\mathbb{R}}
\DeclareMathOperator{\Tr}{Tr}

% ---- Cajas de apoyo --------------------------------------------------------
% "concepto": definiciones previas que hacen falta para entender el capitulo.
\newtcolorbox{concepto}[1]{breakable,colback=blue!4,colframe=blue!55!black,
  fonttitle=\bfseries,title={Concepto previo: #1}}
% "nota": aclaraciones y advertencias.
\newtcolorbox{nota}[1]{breakable,colback=orange!6,colframe=orange!70!black,
  fonttitle=\bfseries,title={#1}}

% ---- Estilos de listados ----------------------------------------------------
\lstdefinestyle{py}{language=Python,basicstyle=\ttfamily\footnotesize,
  breaklines=true,postbreak=\mbox{\textcolor{gray}{$\hookrightarrow$}\space},
  numbers=left,numberstyle=\tiny\color{gray},frame=single,
  commentstyle=\color{green!40!black}\itshape,
  keywordstyle=\color{blue!70!black}\bfseries,
  stringstyle=\color{red!60!black},showstringspaces=false,tabsize=4,
  columns=fullflexible,keepspaces=true}
\lstdefinestyle{salida}{basicstyle=\ttfamily\footnotesize,breaklines=true,
  frame=leftline,backgroundcolor=\color{gray!8},columns=fullflexible,
  keepspaces=true}

  \begin{document}
  \title{\tituloconcepto}
  \author{\autor}
  \date{\today}
  \maketitle
\fi


\section{La esfera de Bloch: un qubit como punto de una esfera}

\subsection{Objetivo}
Entender que todo estado de un qubit se describe con solo dos ángulos,
$\theta$ y $\varphi$, y que por eso se puede dibujar como un punto sobre una
esfera.

\subsection{Conceptos previos}
\begin{concepto}{amplitudes, fase y regla de Born}
Ver capítulo~01 y la sección~\ref{sec:prerreq}. Hace falta saber que un qubit
es $\alpha\ket0+\beta\ket1$ con $|\alpha|^2+|\beta|^2=1$, y que $|\alpha|^2$ y
$|\beta|^2$ son las probabilidades de medir 0 y 1.
\end{concepto}

\subsection{Fundamento teórico}
Los coeficientes de un qubit son números complejos, así que cada uno tiene un
módulo y una fase:
\begin{equation}
  \ket\psi=r_1e^{i\gamma_1}\ket0+r_2e^{i\gamma_2}\ket1 .
\end{equation}
Para simplificar esta expresión se usan dos ideas.

\begin{enumerate}
  \item \textbf{Solo importa la fase relativa.} Multiplicar todo el estado por
    $e^{i\gamma}$ no cambia nada observable (la \emph{fase global} no se puede
    medir). Se puede entonces sacar $e^{i\gamma_1}$ como factor común y
    quedarse solo con la fase relativa $\varphi=\gamma_2-\gamma_1$.
  \item \textbf{Normalización.} Como $r_1^2+r_2^2=1$ se pueden escribir
    $r_1=\cos\frac\theta2$ y $r_2=\sin\frac\theta2$.
\end{enumerate}
El resultado es
\begin{equation}
  \ket\psi=\cos\tfrac\theta2\,\ket0+e^{i\varphi}\sin\tfrac\theta2\,\ket1 .
\end{equation}
Cada qubit queda determinado por dos ángulos, y por tanto es un punto de una
esfera de radio 1: la \textbf{esfera de Bloch}, propuesta por Felix Bloch en
1946. El polo norte es $\ket0$, el polo sur es $\ket1$, el ángulo $\theta$ se
mide desde el eje $z$ y el ángulo $\varphi$ gira alrededor de él. En un qubit
físico de espín, $\theta$ es el ángulo entre el momento angular de espín y el
eje $z$, y $\varphi$ es la fase del espín.

Las coordenadas del punto son los promedios de las matrices de Pauli:
\begin{equation}
  (x,y,z)=\big(\langle X\rangle,\langle Y\rangle,\langle Z\rangle\big)
  =(\sin\theta\cos\varphi,\ \sin\theta\sin\varphi,\ \cos\theta).
\end{equation}
Este mismo resultado se usa en el capítulo~13 (tomografía).

\begin{center}
\begin{tabular}{@{}lll@{}}
\toprule
Estado & Punto $(x,y,z)$ & Dónde está\\
\midrule
$\ket0$ & $(0,0,+1)$ & polo norte\\
$\ket1$ & $(0,0,-1)$ & polo sur\\
$\ket+=(\ket0+\ket1)/\sqrt2$ & $(+1,0,0)$ & ecuador\\
$\ket-=(\ket0-\ket1)/\sqrt2$ & $(-1,0,0)$ & ecuador\\
$(\ket0+i\ket1)/\sqrt2$ & $(0,+1,0)$ & ecuador\\
\bottomrule
\end{tabular}
\end{center}

\paragraph{Ángulo y probabilidades.} Las probabilidades dependen solo de
$\theta$: $P(0)=\cos^2\frac\theta2$ y $P(1)=\sin^2\frac\theta2$. Cambiar
$\varphi$ mueve el punto alrededor del eje $z$ sin tocar $z$, así que no cambia
las probabilidades de esta medición; la fase relativa solo se nota después,
por interferencia (capítulo~05a).

\subsection{Código (completamente comentado)}
\lstinputlisting[style=py]{../../fuentes/10_esfera_bloch.py}

\subsection{Explicación celda por celda}
\begin{description}
  \item[Celda 1] \texttt{Pauli} y \texttt{expectation\_value} calculan
    $\langle X\rangle,\langle Y\rangle,\langle Z\rangle$; con ellos
    \texttt{vector\_bloch} da el punto de la esfera.
  \item[Celda 2] Crea un qubit con $\theta=\pi/3$ y $\varphi=\pi/4$ a partir de
    la fórmula del capítulo.
  \item[Celda 3] Multiplica el estado por $e^{i\gamma}$ y comprueba que el
    punto es idéntico.
  \item[Celda 4] Calcula el punto de cinco estados especiales.
  \item[Celda 5] Compara el punto calculado con la fórmula con $\theta$ y
    $\varphi$.
  \item[Celda 6] Cambia $\varphi$ y muestra que $z$ y $P(1)$ no cambian.
  \item[Celda 7] Dibuja la esfera con \texttt{plot\_bloch\_multivector}.
\end{description}

\subsection{Salida obtenida}
\lstinputlisting[style=salida]{../../fuentes/salidas/10_esfera_bloch.txt}
La salida es determinista (no hay medición).

\subsection{Interpretación}
El punto calculado coincide con la fórmula, con o sin fase global. Al pasar
$\varphi$ de $0$ a $\pi/2$ y a $\pi$, el punto gira de $(0.866,0,0.5)$ a
$(0,0.866,0.5)$ y a $(-0.866,0,0.5)$, pero $P(1)=0.25$ no cambia.

\subsection{Ejercicios}
\begin{enumerate}
  \item ¿Qué $\theta$ da $P(1)=0.9$? Compruébalo con el código.
  \item Calcula a mano el punto de $\ket{-i}=(\ket0-i\ket1)/\sqrt2$ y
    verifícalo con \texttt{Statevector.from\_label("l")}.
  \item ¿En qué lugar de la esfera quedan dos estados que solo difieren en una
    fase global?
\end{enumerate}

\subsection{Errores comunes}
\begin{itemize}
  \item Confundir $\theta$ con $\theta/2$: en la fórmula del estado aparece
    $\cos\frac\theta2$, pero el ángulo que se mide sobre la esfera es
    $\theta$. Por eso $\ket0$ y $\ket1$, que son ortogonales, quedan en polos
    opuestos (a $180^\circ$).
  \item La esfera solo sirve para \emph{un} qubit. Un estado entrelazado de
    varios qubits no se puede separar en puntos de esferas (capítulos~12
    y~13).
\end{itemize}

\ifdefined\enmain\else
  \end{document}
\fi

